A rhombus has perpendicular diagonals that bisect one another. One diagonal is $AC$, of length $\sqrt2$, so the other, $XY$, lies on the perpendicular line through the square centre. Since $X,Y$ are on the incircle, $XY$ is its diameter, of length $1$. The rhombus area is $\dfrac{1}{2} AC\cdot XY=\boxed{\dfrac{\sqrt2}{2}}$.
\begin{asy}
texpreamble("\usepackage{amsmath,amssymb}");
size(240);
defaultpen(fontsize(9pt));
pair A=(0,0),B=(1,0),C=(1,1),D=(0,1),O=(.5,.5),X=O+.5*dir(135),Y=O+.5*dir(-45);draw(A--B--C--D--cycle);draw(circle(O,.5));filldraw(A--X--C--Y--cycle,lightblue);draw(A--C,dashed);draw(X--Y,dashed);label("A",A,SW);label("C",C,NE);label("X",X,NW);label("Y",Y,SE);
\end{asy}
